\section*{Chapitre \ref{ch:g12:metamorphic-rocks} --- Les roches métamorphiques}

\begin{solution}{exo:g12:metamorphic-rocks:1}
(a) des \omterm{def:g12:metamorphic-rocks:polymorph}{polymorphes}~; (b) une \omterm{def:g12:metamorphic-rocks:index-mineral}{isograde}~; (c) une \omterm{def:g12:metamorphic-rocks:contact-regional}{auréole}~; (d) une
\omterm{def:g12:metamorphic-rocks:foliation}{foliation}~; (e) la \omterm{def:g12:metamorphic-rocks:lithostatic}{pression lithostatique}.
\end{solution}

\begin{solution}{exo:g12:metamorphic-rocks:2}
(1) un marbre~; (2) un quartzite~; (3) une ardoise~; (4) un gneiss~; (5) une
cornéenne.
\end{solution}

\begin{solution}{exo:g12:metamorphic-rocks:3}
$P = \rho g z$~: à \qty{10}{km}, $2830 \times 9.81 \times \num{10000} \approx
\qty{2.78e8}{Pa} = \qty{0.278}{GPa}$~; à \qty{20}{km}, deux fois plus,
\qty{0.555}{GPa}.
\end{solution}

\begin{solution}{exo:g12:metamorphic-rocks:4}
(1) Le disthène (la limite disthène--andalousite est vers
\qty{0.25}{GPa} à \qty{400}{\celsius}). (2) L'andalousite. (3) La
sillimanite.
\end{solution}

\begin{solution}{exo:g12:metamorphic-rocks:5}
(1) Ardoise, phyllade, schiste, gneiss. (2) Dans le schiste. (3) Le gneiss.
\end{solution}

\begin{solution}{exo:g12:metamorphic-rocks:6}
(1) $2830 \times 9.81 \times \num{35000} \approx \qty{9.72e8}{Pa}$, soit
$\num{9.72e8}/101325 \approx \num{9600}$ atmosphères. (2) $1025 \times 9.81
\times 4000 \approx \qty{4.02e7}{Pa} = \qty{40.2}{MPa}$. (3) $972/40.2
\approx 24$ fois.
\end{solution}

\begin{solution}{exo:g12:metamorphic-rocks:7}
(1) $z = \num{0.375e9}/(2830 \times 9.81) \approx \qty{13.5}{km}$.
(2) $(504 - 10)/13.5 \approx \qty{36.6}{\celsius}$ par kilomètre. (3) Non~:
la roche devrait se trouver exactement en un point du diagramme~; de telles
roches sont rares.
\end{solution}

\begin{solution}{exo:g12:metamorphic-rocks:8}
(1) $23/3 \approx \qty{8}{mm}$. (2) Environ 8 millions d'années. (3) La
roche est restée en profondeur, à haute température, pendant des millions
d'années.
\end{solution}

\begin{solution}{exo:g12:metamorphic-rocks:9}
(1) Quatre~: les \omterm{def:g12:metamorphic-rocks:index-mineral}{isogrades} du grenat, de la staurotide, du disthène et de
la sillimanite, dans cet ordre. (2) Vers l'ouest. (3) À l'extrémité ouest,
dans la zone à sillimanite. Ses roches ont franchi la limite
disthène--sillimanite~: enfouies à environ \qty{20}{km}
(\qty{0.56}{GPa}), elles étaient plus chaudes qu'environ
\qty{590}{\celsius}.
\end{solution}

\begin{solution}{exo:g12:metamorphic-rocks:10}
(1) La boue devient une cornéenne, le calcaire un marbre, le grès un
quartzite. (2) Elles ont été chauffées sans être comprimées~: aucune
pression orientée n'a aligné leurs minéraux. (3) Les roches qui l'entourent
sont déjà chaudes~: le magma ne les chauffe guère davantage.
\end{solution}

\begin{solution}{exo:g12:metamorphic-rocks:11}
(1) En travers de la poussée~: leurs faces plates sont tournées vers l'est
et l'ouest, si bien que les paillettes se dressent dans des plans
nord--sud. (2) Verticale, orientée nord--sud. (3) Les grains de calcite et
de quartz ne sont ni plats ni allongés~: même quand ils recristallisent
sous une compression, il n'y a rien à aligner.
\end{solution}

\begin{solution}{exo:g12:metamorphic-rocks:12}
(1) $z = (T - 10)/30$, donc $P = 0.0278\,(T - 10)/30 \approx
0.000927\,(T - 10)$. (2) $0.000927\,(T - 10) = 0.577 + 0.00209\,(T - 600)$
donne $0.001163\,T \approx 0.668$, donc $T \approx \qty{574}{\celsius}$, $P
\approx 0.000927 \times 564 \approx \qty{0.523}{GPa}$ et $z = (574 - 10)/30
\approx \qty{18.8}{km}$. (3) Le disthène au-dessus de cette profondeur, la
sillimanite au-dessous.
\end{solution}

\begin{solution}{exo:g12:metamorphic-rocks:13}
(1) $10 + 10 \times 35 = \qty{360}{\celsius}$. (2) $10 + 30 \times 35 =
\qty{1060}{\celsius}$. (3) La pression à \qty{35}{km} vaut environ
\qty{0.97}{GPa}~; à \qty{360}{\celsius}, ce point se trouve au cœur du
domaine du disthène. Le second point se trouve loin à droite du diagramme,
dans le domaine de la sillimanite, à des températures où les roches de la
croûte commencent à fondre. (4) Il leur faut une haute pression à basse
température, une combinaison que seule fournit la croûte froide entraînée
vers le bas dans une \omterm{def:g11:subduction-zones:subduction}{zone de subduction}.
\end{solution}

\begin{solution}{exo:g12:metamorphic-rocks:14}
(1) Disthène~: $1000/3.61 \approx \qty{277}{cm^3}$~; sillimanite~: $1000/3.24
\approx \qty{309}{cm^3}$~; andalousite~: $1000/3.15 \approx \qty{317}{cm^3}$.
(2) $(317 - 277)/317 \approx 0.127$, soit environ \qty{13}{\percent}.
(3) À haute pression, l'arrangement le plus compact des atomes est
favorisé~; le disthène loge les mêmes atomes dans le moins de place.
\end{solution}

\begin{solution}{exo:g12:metamorphic-rocks:15}
(1) \qty{0.555}{GPa} à \qty{20}{km}, \qty{0.694}{GPa} à \qty{25}{km}.
(2) $T = 600 + (P - 0.577)/0.00209$~: environ \qty{590}{\celsius} à
\qty{0.555}{GPa} et \qty{656}{\celsius} à \qty{0.694}{GPa}. À
\qty{700}{\celsius}, les roches étaient plus chaudes que ces deux valeurs,
dans le domaine de la sillimanite~: c'est compatible. (3) $(700 - 10)/20 \approx \qty{34.5}{\celsius}$
et $(700 - 10)/25 \approx \qty{27.6}{\celsius}$ par kilomètre, proches du
gradient typique. (4) De $\qty{20}{km}$ à $\qty{25}{km}$ en 40 millions
d'années~: de 0.50 à \qty{0.63}{km} par million d'années, soit de 0.50 à
\qty{0.63}{mm} par an.
\end{solution}

\begin{solution}{pb:g12:metamorphic-rocks:1}
\textbf{1.} Elle est foliée, et ses micas sont assez grands pour être vus.
Ce sont les micas, la muscovite et la biotite, qui forment la \omterm{def:g12:metamorphic-rocks:foliation}{schistosité}.

\textbf{2.} \ce{Al2SiO5}~; ce sont des \omterm{def:g12:metamorphic-rocks:polymorph}{polymorphes}.

\textbf{3.} Les volumes d'un kilogramme sont dans le rapport
$3.15/3.61 \approx
0.873$~: le disthène est plus petit d'environ
\qty{13}{\percent}. La pression favorise le disthène.

\textbf{4.} La zone à disthène. Pendant sa descente, elle a franchi les
\omterm{def:g12:metamorphic-rocks:index-mineral}{isogrades} de la biotite, du grenat, de la staurotide et du disthène.

\textbf{5.} Environ 8 millions d'années.

\textbf{6.} $P = 0.577 + 0.00209 \times 0 = \qty{0.577}{GPa}$.

\textbf{7.} \qty{0.70}{GPa} est au-dessus de \qty{0.577}{GPa}, du côté du
disthène de la limite disthène--sillimanite~; c'est aussi bien au-dessus du
domaine de l'andalousite, qui se trouve à \qty{600}{\celsius} sous environ
\qty{0.23}{GPa}.

\textbf{8.} \qty{0.577}{GPa}.

\textbf{9.} $0.70 = 0.577 + 0.00209\,(T - 600)$, donc $T = 600 +
0.123/0.00209 \approx \qty{659}{\celsius}$.

\textbf{10.} \qty{7.0}{kbar}~; $\num{7.0e8}/101325 \approx \num{6900}$
atmosphères.

\textbf{11.} $\rho g = 2830 \times 9.81 \approx \qty{27800}{Pa}$ par mètre,
soit \qty{27.8}{MPa} par kilomètre.

\textbf{12.} $z = \num{0.70e9}/\num{27762} \approx \qty{25200}{m}$, environ
\qty{25.2}{km}.

\textbf{13.} $\num{0.577e9}/\num{27762} \approx \qty{20.8}{km}$.

\textbf{14.} $\num{0.70e9}/(2700 \times 9.81) \approx \qty{26.4}{km}$, soit
environ \qty{5}{\percent} plus profond ($26.4/25.2 \approx 1.048$).

\textbf{15.} Dans la croûte~: \qty{25}{km}, c'est moins que l'épaisseur
même d'une croûte normale.

\textbf{16.} $(600 - 10)/25.2 \approx \qty{23.4}{\celsius}$ par kilomètre.

\textbf{17.} Il est bien plus chaud que celui d'une \omterm{def:g11:subduction-zones:subduction}{zone de subduction} (moins
de \qty{10}{\celsius} par kilomètre) et un peu plus froid que celui d'une
croûte typique (\qty{30}{\celsius}), loin de celui d'une région volcanique
(40 à \qty{50}{\celsius})~: c'est une roche enfouie profondément dans la
croûte épaissie d'une chaîne de montagnes, par \omterm{def:g12:metamorphic-rocks:contact-regional}{métamorphisme régional}.

\textbf{18.} $10 + 30 \times 25.2 \approx \qty{766}{\celsius}$. La limite
disthène--sillimanite, prolongée, se trouve à $0.577 + 0.00209 \times 166
\approx \qty{0.92}{GPa}$ à cette température~; \qty{0.70}{GPa} est
au-dessous, dans le domaine de la sillimanite. Un \omterm{def:g11:lithosphere-asthenosphere:geotherm}{géotherme} typique aurait
produit de la sillimanite.

\textbf{19.} $\qty{25.2}{km} = \qty{2.52e7}{mm}$~; à \qty{0.5}{mm} par an,
environ 50 millions d'années.

\textbf{20.} En remontant, la roche s'est refroidie, si bien que ses
réactions ont ralenti, et l'eau dont les réactions inverses ont besoin
avait été chassée pendant la descente.

\textbf{21.} Le schiste à disthène a cristallisé à environ
\textbf{\qty{25}{km} de profondeur}, vers \qty{600}{\celsius}.
\end{solution}
